\documentclass{report}

% Page Size
\usepackage[letterpaper, portrait, margin=1in, left=1.5in]{geometry}

% Spacing
\usepackage{setspace}
\setlength{\parskip}{\baselineskip}
\doublespacing{}


% Bibliography
\usepackage[american]{babel}
\usepackage[style=apa, citestyle=apa, backend=biber]{biblatex}
\DeclareLanguageMapping{american}{american-apa}
\addbibresource{references.bib}
\usepackage[babel,threshold=2]{csquotes}

% Graphics
\usepackage{graphicx}
\graphicspath{ {./images/} }

% Table of Content
\setcounter{tocdepth}{2}

% Code Formatting
\usepackage{minted}
\usemintedstyle{pastie}

\def\code#1{\texttt{#1}}

% Positioning
\usepackage{float}

% Colors
\usepackage[dvipsnames]{xcolor}

% Links
\usepackage{hyperref}
\hypersetup{
	pdftitle={Development of a touch typing trainer with an emphasis on finger and wrist placement},
	pdfstartview={FitH},
	colorlinks=true,
	linkcolor=black,
	citecolor=black,
	filecolor=black,
	urlcolor=MidnightBlue
}


\begin{document}

\pagenumbering{gobble}
\begin{titlepage}
	\centering
	\includegraphics[width=0.25\textwidth]{upc.png}
	\includegraphics[width=0.25\textwidth]{dcs.png}
	\par\vspace{1cm}
	{\scshape\LARGE Research Brief \par}
	\par\vspace{1cm}
	\begin{tabular}{ | l | l | }
		\hline
		Name               & Oscar Vian L. Valles                  \\
		\hline
		Student ID number  & 2018--06491                           \\
		\hline
		Project Title      & Development of a touch typing trainer \\ & with an emphasis on finger and wrist positions\\
		\hline
		Project No.        & 9                                     \\
		\hline
		Name of Supervisor & Dhong Fhel K. Gom-os                  \\
		\hline
	\end{tabular}




\end{titlepage}
\newpage

\tableofcontents
\newpage

\pagenumbering{arabic}

\setcounter{chapter}{1}
\section{Background of the Study}
There are a lot of educational typing tests available that help people learn
touch typing, including Monkeytype, TypeRacer, or Keybr. These typing tests list
out words that are then typed out. The inputted keys are then compared to check
if the user has typed the expected letter. At the end of the test, the time
taken is calculated and certain metrics is given. These metrics include words
per minute (WPM) and accuracy \parencite{bartnik2021}.

However, this method of examination leaves out a crucial part of typing —
ergonomics. Ergonomic typing prevents a lot health issues in the future like
repetitive strain injury or carpal tunnel. One important factor that affects
ergonomics is the typing procedure and posture. This means proper placement of
the wrist, hands, and hitting the keys using the right finger that is assigned
to the key.

Correct finger placement is usually taught at the beginning using a diagram,
with each key being associated with a specific finger. For instance, the letter
Q in a QWERTY layout should be hit using the fifth digit of the left hand, and
this is shown by coloring the fifth digit and the key Q with the same color or
by placing the letters directly on the fingers \parencite{dobson2009touch}.

Incorrect finger placement may cause these hand and wrist positions: ulnar
deviation, forearm pronation, and wrist extension \parencite{serina1999}. These
three are hand and wrist positions that are common in all activities, however,
prolonged periods in these positions may cause cause injuries such as Carpal
tunnel syndrome (CTS) \parencite{toosi2015}

In addition, this type of typing is frequently taught in the beginner level.
However, finger position is not \parencite{donica2018}

There is currently no way to automatically detect which finger is used during
typing. No hardware or software solutions currently exist in the market. This
research aims to solve this by combine finger tracking, with educational typing
tests to determine which finger pressed the key and if that finger is the
correct one to use.

\section{Significance of the Research}
This research is beneficial for all users of physical keyboards. These include a
vast majority of the population as there are a lot of professions that heavily
rely on keyboards. Examples include developers, physicians, educators,
accountants.

This research has a direct impact on people that has onset RSI or other
hand/wrist injuries [ statistics ] that are caused by poor typing habits. By
correcting these poor habits, pain from these injuries will be lessened, and
even be prevented from occurring in the first place.

\newpage
\nocite{*} \printbibliography[heading=bibintoc,title={References}]{}

\end{document}
